% ============================================================
%  CHAPTER 3: SYSTEM DESIGN AND ARCHITECTURE
%  AgriVerify — AI-Powered Fake Seed Detection System
%  Manipur Technical University, Final Year Project Report
% ============================================================

\chapter{System Design and Architecture}

% ──────────────────────────────────────────────────────────────
\section{Overall System Architecture}
% ──────────────────────────────────────────────────────────────

AgriVerify is engineered as a distributed, multi-tier web architecture that seamlessly integrates deep learning model inference, cloud-based AI orchestration, and a role-aware web application into a unified platform. The system adheres to \textbf{Clean Architecture} principles (specifically the Onion Architecture variant) as described by Martin \cite{martin2008clean}, ensuring complete separation between core business logic and infrastructure concerns such as databases, cloud vendor dependencies, and user interface frameworks. This architectural approach ensures that dependencies flow strictly inward: outer layers (Infrastructure, API) depend exclusively on abstractions exposed by inner layers (Application, Domain), never the reverse. This dependency inversion principle promotes maintainability, testability, and adaptability to future technology changes.

\medskip
The system comprises three major subsystems:

\begin{enumerate}
  \item \textbf{Deep Learning Microservice} — A PyTorch-based ResNet-50 convolutional neural network \cite{he2016deep} trained on a corpus of 1{,}563 paddy seed images, served as a containerized FastAPI endpoint on Google Cloud Platform (GCP) Cloud Run, enabling real-time seed classification and physical parameter extraction.

  \item \textbf{Backend REST API} — An ASP.NET Core 8 Web API \cite{microsoft2023aspnetcore} functioning as the primary orchestration layer. It manages user authentication, mediates access to the ML microservice and Azure AI services (Custom Vision, AI Vision OCR, OpenAI), persists results to Microsoft SQL Server through Entity Framework Core, and exposes a versioned REST API consumed by the frontend application.

  \item \textbf{Frontend Web Application} — A Next.js 14 App Router application \cite{nextjs2024} providing role-based dashboards for Farmers, Agricultural Officers, and Administrators. All browser-to-backend communication is forwarded through a server-side proxy route to eliminate credential exposure in the client environment and centralize authentication token management.
\end{enumerate}

\medskip
The four architectural layers of the backend system follow the Onion Architecture pattern:

\begin{enumerate}
  \item[Domain Layer] Contains core domain entities (\texttt{User}, \texttt{VerificationHistory}, \texttt{ProductComplaint}), value objects, domain-specific enumerations (\texttt{UserRole}, \texttt{VerificationStatus}, \texttt{ComplaintStatus}), and domain-specific exceptions. This layer is completely framework-agnostic and carries no dependencies on external libraries or services, ensuring pure domain logic isolation.

  \item[Application Layer] Orchestrates business workflows through application services (\texttt{VerificationService}, \texttt{UserService}, \texttt{ComplaintService}), defines Data Transfer Objects (DTOs) for inter-layer communication, configures AutoMapper profiles for entity mapping, and implements FluentValidation rules for input validation. This layer depends exclusively on the Domain layer, maintaining high cohesion and low coupling.

  \item[Infrastructure Layer] Provides concrete implementations for interfaces declared in the Application layer: Entity Framework Core repositories, the \texttt{ApplicationDbContext} database context, Azure service wrappers (\texttt{CustomVisionService}, \texttt{ComputerVisionService}, \texttt{OpenAIService}, \texttt{BlobStorageService}), and database migration management. This layer acts as a bridge between the application and external systems without exposing implementation details to upper layers.

  \item[API Layer] Hosts ASP.NET Core controllers (\texttt{AuthController}, \texttt{VerificationController}, \texttt{ComplaintController}, \texttt{AdminController}), the JWT middleware authentication pipeline, Swagger/OpenAPI 3.0 documentation, and comprehensive global exception handling middleware for consistent error response formatting.
\end{enumerate}


% ──────────────────────────────────────────────────────────────
\section{Functional Requirements}
% ──────────────────────────────────────────────────────────────

The following functional requirements were systematically derived from stakeholder analysis and constitute the foundation of the system implementation:

\subsection*{FR-01: User Registration and Authentication}

\begin{itemize}
  \item Farmers must register with their name, email address, geographic location (state and district), and authenticate using email and password credentials.
  \item Agricultural officers must authenticate using their officer ID and a time-based one-time password (TOTP) as a second authentication factor, enhancing security for privileged operations.
  \item System administrators are provisioned through database seeding with dedicated credentials managed by the platform's configuration system.
  \item Upon successful authentication, the system must issue JWT access tokens (short-lived) and refresh tokens (long-lived) to enable stateless session management.
  \item Session invalidation must occur immediately upon user logout, with token revocation mechanisms in place to prevent replay attacks.
\end{itemize}

\subsection*{FR-02: Seed Verification}

\begin{itemize}
  \item Authenticated farmers must upload front-facing and back-facing photographs of seed packaging for AI-powered authenticity analysis.
  \item The system must execute parallel inference tasks using Azure Custom Vision for package classification, Azure AI Vision for optical character recognition (OCR) text extraction, and Azure OpenAI for contextual report synthesis.
  \item The verification result must include a machine-learning confidence score, categorical verification status (Genuine / Counterfeit / Suspicious), extracted textual information, 17 dimensionality-uniformity-stability (DUS) physical parameters, and a natural-language summary.
  \item An unauthenticated public verification endpoint must remain available to accommodate anonymous users seeking verification without account creation.
  \item All verification records must be persisted with complete audit trails, including submission timestamps, user identifiers, and result parameters.
\end{itemize}

\subsection*{FR-03: Complaint Management}

\begin{itemize}
  \item Authenticated farmers must file formal complaints against suspicious seed batches, specifying issue type (germination failure, phenotypic abnormality, packaging defect), geographic district, severity level (low, medium, high), and optional photographic evidence.
  \item Agricultural officers must view complaints assigned to their jurisdiction, update complaint status through a workflow (Open → Investigation → Resolved → Closed), record investigation actions with automated timestamps, and attach evidentiary attachments.
  \item System administrators must have visibility into all complaints across all geographic regions, filter by district and status, assign unassigned complaints to specific officers, and access aggregated complaint statistics for policy decisions.
\end{itemize}

\subsection*{FR-04: Role-Based Dashboards}

\begin{itemize}
  \item Farmer dashboards must display personal verification history, current complaint statuses, summary metrics (total verifications, genuine vs. counterfeit distribution), and recommended actions.
  \item Officer dashboards must display complaint queue prioritized by severity and date, enable batch risk assessment across multiple complaints, provide district-level analytics, and support complaint reassignment.
  \item Administrator dashboards must provide platform-wide analytics, user management capabilities, reference data configuration, system health monitoring, and export functionality for compliance reporting.
\end{itemize}

\subsection*{FR-05: Batch Risk Registry}

\begin{itemize}
  \item The system must maintain a \texttt{BatchRiskRegistry} entity that aggregates complaint statistics per seed batch number, tracks geographic distribution of affected districts, computes average severity scores, assigns dynamic risk levels (low, medium, high, critical) based on complaint volume and severity patterns.
  \item Risk registry updates must be triggered automatically on every complaint state transition, maintaining real-time accuracy for officer-facing risk assessments.
\end{itemize}

\subsection*{FR-06: Deep Learning Model Inference}

\begin{itemize}
  \item The ResNet-50 microservice must accept seed images, apply the standardized 224\,×\,224-pixel preprocessing pipeline with ImageNet normalization \cite{he2016deep}, execute forward neural network inference, and return class probability distributions.
  \item The microservice must extract and return 17 dimensionality-uniformity-stability (DUS) physical parameters (grain length, grain width, grain thickness, color distribution statistics, texture metrics, shape metrics) using OpenCV morphological analysis \cite{opencv2020}.
  \item Variety matching must occur against a reference database of 12 registered paddy varieties, returning the top-3 highest-similarity matches with quantified similarity percentages.
\end{itemize}

% ──────────────────────────────────────────────────────────────
\section{Non-Functional Requirements}
% ──────────────────────────────────────────────────────────────

\begin{description}
  \item[NFR-01 — Performance]
        The end-to-end verification pipeline (image upload → ML inference → database persistence → API response) must complete within interactive latency bounds acceptable for web-based applications. Parallel execution of Azure OCR and Custom Vision operations via \texttt{Task.WhenAll} orchestration is mandated to minimize per-request latency through concurrent I/O. The FastAPI microservice must serve single-image inference in sub-second wall-clock time on GCP Cloud Run hardware \cite{fastapi2024}, and cumulative API response time must not exceed 5 seconds under normal load conditions.

  \item[NFR-02 — Scalability]
        The backend must be designed as a stateless distributed system deployable as Docker containers across multiple computational instances, enabling horizontal scaling without code modifications. The frontend Next.js application must support server-side rendering (SSR) for initial page loads and leverage edge caching through Content Delivery Networks (CDN) to reduce perceived latency for geographically distributed users \cite{nextjs2024}.

  \item[NFR-03 — Security]
        All HTTP traffic must be transmitted exclusively over Transport Layer Security (TLS) version 1.3 or higher, enforcing HTTPS on all endpoints. JWT access tokens must carry minimal claim sets and expire within a configurable time-to-live (TTL), typically 15–30 minutes. User passwords must be stored as cryptographic hashes using PBKDF2 with adequate iteration counts or bcrypt with appropriate cost factors. Sensitive credentials (Azure subscription keys, database connection strings, API authentication tokens) must be injected through environment variables and never embedded in source code repositories. Role-based endpoint protection must be enforced through ASP.NET Core \texttt{[Authorize]} attributes with policy-based claim validation \cite{microsoft2023aspnetcore}.

  \item[NFR-04 — Reliability and Data Integrity]
        All grouped repository operations must execute within ACID-compliant database transactions with atomic commit-or-rollback semantics, ensuring consistent system state across distributed operations. Soft-delete functionality via \texttt{IsDeleted} and \texttt{DeletedAt} timestamps must be implemented for complaint and verification records to preserve historical audit trails for future legal compliance requirements. Health check endpoints (\texttt{/health}, \texttt{/api/health}) must expose real-time connectivity status of critical dependencies (SQL Server, Azure services, ML microservice).

  \item[NFR-05 — Maintainability]
        The codebase must rigorously adhere to Clean Architecture layering patterns with explicit separation of concerns between Domain, Application, Infrastructure, and API layers. Infrastructure services must be registered through dedicated \texttt{InfrastructureServiceCollectionExtensions} extension classes rather than inline service registration in \texttt{Program.cs}, promoting code organization and testability. All public API endpoints must be documented through Swagger/OpenAPI 3.0 specification, enabling automated API client generation and reducing integration friction.

  \item[NFR-06 — Usability and Accessibility]
        The frontend must deliver a responsive layout functional on both desktop browsers and mobile devices, with particular attention to farmers accessing the platform via personal smartphones with variable connectivity. The user interface must provide immediate loading feedback through skeleton loading placeholders and non-blocking toast notifications for asynchronous operations, reducing user uncertainty during data submission.

  \item[NFR-07 — Portability]
        The deep learning model must be exported in multiple serialization formats (PyTorch native \texttt{.pth}, TorchScript \texttt{.pt}, Open Neural Network Exchange ONNX \texttt{.onnx}) to support deployment across heterogeneous runtime environments. The entire application stack (frontend, backend, ML microservice) must be containerizable through Docker, enabling consistent deployment across development, staging, and production environments.
\end{description}

% ──────────────────────────────────────────────────────────────
\section{Technology Stack and Tools}
% ──────────────────────────────────────────────────────────────

\subsection{Deep Learning Stack}

The seed classification model is implemented using \textbf{Python 3.10} with the following specialized libraries:

\begin{itemize}
  \item \textbf{PyTorch and TorchVision} — PyTorch \cite{pytorch2024} serves as the primary deep learning framework, providing dynamic computational graphs that simplify debugging and enable flexible architecture modifications during transfer-learning fine-tuning phases. TorchVision provides pre-trained model architectures and image preprocessing utilities.

  \item \textbf{ResNet-50 Architecture} — ResNet-50 \cite{he2016deep} is a 50-layer deep residual network pre-trained on the ImageNet dataset (1.2 million images, 1{,}000 object classes). The backbone contributes approximately 23.5 million parameters; a custom three-layer classification head (2{,}048 → 512 → 128 → 3 units) adds an additional 1.1 million trainable parameters optimized for the seed classification task.

  \item \textbf{OpenCV} — The OpenCV library \cite{opencv2020} enables the dimensionality-uniformity-stability (DUS) parameter extraction pipeline, providing robust contour detection, morphological transformations, and pixel-to-millimetre calibration for quantifying seed physical characteristics.

  \item \textbf{Numerical Computing} — NumPy and Pandas \cite{numpy2024} enable efficient numerical array operations, statistical calculations, and dataset manipulation. scikit-learn provides stratified sampling for train-validation-test splitting with statistical validity \cite{scikit2024}.

  \item \textbf{Visualization} — Matplotlib and Seaborn enable training curve visualization and confusion matrix rendering for model performance analysis and debugging.
\end{itemize}

\noindent
Model training employed the AdamW optimizer with weight decay ($\lambda = 10^{-4}$), cross-entropy loss with label smoothing ($\varepsilon = 0.1$), and a two-phase transfer learning schedule. Phase 1 (Epochs 1–8) froze the backbone weights with learning rate $10^{-3}$ to stabilize the classification head. Phase 2 (Epochs 9–15) unfroze all weights with learning rate $10^{-4}$ for fine-tuning backbone representations. The final ResNet-50 model achieved 99.57\% test accuracy on the combined 1{,}563-image dataset with a train-validation gap below 5\%, demonstrating exceptional generalization capability.

% ──────────────────────────
\subsection{Backend Technology Stack}

The backend REST API is constructed on the \textbf{Microsoft .NET 8} ecosystem. Table~\ref{tab:backend_stack} summarizes the key architectural components:

\begin{table}[h!]
  \centering
  \caption{Backend Technology Stack Components}
  \label{tab:backend_stack}
  \renewcommand{\arraystretch}{1.5}
  \begin{tabularx}{\textwidth}{p{3.5cm}X}
    \hline
    \textbf{Component} & \textbf{Technology} \\
    \hline
    Web Framework & ASP.NET Core 8.0 / .NET 8 (Kestrel HTTP server) \cite{microsoft2023aspnetcore} \\
    \hline
    Database Engine & Microsoft SQL Server 2022 \\
    \hline
    Object-Relational Mapping & Entity Framework Core 8.0 with code-first migrations \\
    \hline
    Authentication & ASP.NET Core Identity + JWT Bearer tokens \cite{microsoft2023aspnetcore} \\
    \hline
    Object Mapping & AutoMapper for DTO transformation \\
    \hline
    Input Validation & FluentValidation + DataAnnotations \cite{fluentvalidation2024} \\
    \hline
    Logging Framework & Serilog with console and file sinks \cite{serilog2024} \\
    \hline
    API Documentation & Swagger / OpenAPI 3.0 with Swashbuckle \\
    \hline
    Cloud AI Services & Azure Custom Vision, Azure AI Vision (OCR), Azure OpenAI \\
    \hline
    Blob Storage & Azure Blob Storage for image persistence \\
    \hline
    Deployment & Docker containerization / Render Cloud platform \\
    \hline
  \end{tabularx}
\end{table}

\noindent
Azure AI services are consumed through dedicated service wrapper classes (\texttt{CustomVisionService}, \texttt{ComputerVisionService}, \texttt{OpenAIService}) registered in the dependency injection container as singletons for resource efficiency. The parallel orchestration pattern using \texttt{Task.WhenAll(ocrTask, classificationTask)} ensures that OCR text extraction and image classification execute concurrently, significantly reducing per-request latency through efficient asynchronous I/O multiplexing \cite{microsoft2023aspnetcore}.

% ──────────────────────────
\subsection{Machine Learning Model Serving}

The trained PyTorch model is deployed as a containerized microservice using \textbf{FastAPI} \cite{fastapi2024} on \textbf{Google Cloud Platform (GCP) Cloud Run}, enabling automatic horizontal scaling based on request volume. The inference pipeline executes through five sequential stages:

\begin{enumerate}
  \item \textbf{Request Reception} — The FastAPI endpoint receives multipart form-data containing the seed image payload.
  \item \textbf{Image Preprocessing} — The image is resized to 224\,×\,224 pixels and normalized using ImageNet channel-wise statistics (mean $\mu = [0.485, 0.456, 0.406]$, standard deviation $\sigma = [0.229, 0.224, 0.225]$) \cite{he2016deep}.
  \item \textbf{Neural Network Inference} — The ResNet-50 model processes the normalized tensor and produces raw class logits.
  \item \textbf{Result Post-processing} — Softmax activation converts logits to probability distributions; OpenCV routines extract 17 DUS parameters; variety matching identifies top-3 reference varieties by Euclidean distance in feature space.
  \item \textbf{JSON Response} — A structured JSON payload containing prediction label, confidence score, physical parameters, DUS classification, and similarity-ranked variety matches is returned to the caller.
\end{enumerate}

\noindent
The model is exported in three formats to maximize deployment flexibility: PyTorch native (\texttt{.pth}, ~98 MB), TorchScript (\texttt{.pt}, ~97 MB for C++ production runtimes), and ONNX (\texttt{.onnx}, ~98 MB for cross-platform inference engines) \cite{pytorch2024}. Continuous integration/continuous deployment (CI/CD) is automated through GitHub Actions, triggering Docker image builds and Artifact Registry pushes on every commit to the main branch.

% ──────────────────────────
\subsection{Frontend Application Stack}

The user-facing application is implemented as a \textbf{Next.js 14 App Router} single-page application written in \textbf{TypeScript}. The technology choices are:

\begin{itemize}
  \item \textbf{Next.js 14 App Router} \cite{nextjs2024} — Provides file-system-based routing with route groups that cleanly partition the three role namespaces (\texttt{(farmer)}, \texttt{(officer)}, \texttt{(admin)}). React Server Components handle initial data fetching for SEO optimization; Client Components manage interactive UI elements.

  \item \textbf{TypeScript} — Enforces strict compile-time type safety throughout the application, including shared Data Transfer Object (DTO) interfaces in \texttt{src/types/}, domain service modules in \texttt{src/services/}, and custom React hooks in \texttt{src/hooks/} \cite{typescript2024}.

  \item \textbf{TanStack Query} — Manages asynchronous server state with intelligent cache invalidation, background refetching, cursor-based pagination, and optimistic mutation rollback \cite{tanstackquery2024}. Query hooks are co-located with domain modules; a global \texttt{QueryClient} configuration disables automatic refetch-on-window-focus to minimize incidental network requests.

  \item \textbf{Zustand} — Provides lightweight client-side state management for authentication context (\texttt{user}, \texttt{authToken}, \texttt{role}, \texttt{isAuthenticated}). The authentication store is initialized from browser localStorage on client mount and exposes \texttt{setAuth} and \texttt{clearAuth} actions.

  \item \textbf{Axios HTTP Client} — Centralizes HTTP request configuration with a base URL targeting the proxy route \texttt{/api/proxy}. Request interceptors automatically inject JWT Bearer tokens from localStorage; response interceptors handle 401 unauthorized responses by transparently requesting token refresh and replaying the original request.

  \item \textbf{Next.js API Route Proxy} — A server-side proxy route at \texttt{/api/proxy/[...path]} forwards browser requests to the ASP.NET backend stored in the \texttt{BACKEND\_API\_URL} environment variable, preventing credential exposure in client-side code while centralizing CORS handling.

  \item \textbf{Tailwind CSS} — Utility-first CSS framework with extended design tokens (color palette, border radii, shadow effects) configured in \texttt{tailwind.config.ts} \cite{tailwindcss2024}.

  \item \textbf{Shadcn/UI Component Library} — Provides production-ready accessible UI components (Button, Input, Select, Dialog, Sheet, Table, Badge, Tabs, Skeleton, Progress) implemented on native HTML5 with minimal abstraction overhead and full TypeScript support.
\end{itemize}

% ──────────────────────────────────────────────────────────────
\section{System Architecture Diagrams}
% ──────────────────────────────────────────────────────────────

\subsection{High-Level Architecture}

Figure~\ref{fig:high_level_arch} illustrates the complete end-to-end system topology. The user's browser communicates exclusively with the Next.js frontend server through HTTPS. All backend API calls are proxied server-side through the Next.js proxy route to the ASP.NET Core API, which orchestrates requests to the ML microservice on GCP and to the suite of Azure AI services. Seed images are persisted in Azure Blob Storage; verification and complaint data are stored in Microsoft SQL Server via Entity Framework Core.

\begin{figure}[h!]
  \centering
  \begin{tikzpicture}[
      node distance=1.8cm and 2.4cm,
      box/.style={rectangle, draw, rounded corners=4pt,
                  minimum width=3.6cm, minimum height=1.0cm,
                  align=center, font=\small},
      cloud/.style={box, fill=blue!8},
      api/.style={box, fill=green!10},
      ml/.style={box, fill=orange!12},
      db/.style={box, fill=gray!12},
      arr/.style={->, >=stealth, thick}
    ]

    % Browser Layer
    \node[box, fill=yellow!15] (browser) {Browser / Mobile Device\\ (Farmer / Officer / Admin)};

    % Frontend Layer
    \node[cloud, below=of browser] (nextjs) {Next.js 14\\ App Router};

    % Backend API Layer
    \node[api, below=of nextjs] (aspnet) {ASP.NET Core 8\\ REST API};

    % Service Layer (horizontal distribution)
    \node[ml, below left=1.8cm and 1.8cm of aspnet]   (fastapi)  {FastAPI Microservice\\ ResNet-50 Classification\\ (GCP Cloud Run)};
    \node[cloud, below=1.8cm of aspnet]               (azure)    {Azure AI Services\\ Custom Vision, OCR,\\ OpenAI Integration};
    \node[db, below right=1.8cm and 1.8cm of aspnet]  (sqlserver){SQL Server 2022\\ + Azure Blob Storage};

    % Connection Arrows
    \draw[arr] (browser) -- node[right, font=\tiny]{HTTPS} (nextjs);
    \draw[arr] (nextjs)  -- node[right, font=\tiny]{Proxy Route} (aspnet);
    \draw[arr] (aspnet)  -- node[above left,  font=\tiny]{REST/JSON} (fastapi);
    \draw[arr] (aspnet)  -- node[right, font=\tiny]{SDK} (azure);
    \draw[arr] (aspnet)  -- node[above right, font=\tiny]{EF Core} (sqlserver);

    % Return arrows (dashed)
    \draw[arr, dashed] (fastapi)  -- (aspnet);
    \draw[arr, dashed] (azure)    -- (aspnet);
    \draw[arr, dashed] (sqlserver)-- (aspnet);
  \end{tikzpicture}
  \caption{High-level system architecture of the AgriVerify platform, showing distributed deployment across browser, frontend, backend, and cloud services.}
  \label{fig:high_level_arch}
\end{figure}

\subsection{Verification Pipeline Data Flow}

Figure~\ref{fig:dfd} illustrates the data flow for the core seed verification workflow. When a farmer submits seed images through the frontend, the request traverses the Next.js proxy route, enters the ASP.NET verification orchestration pipeline, and is distributed in parallel to three independent services: the ML microservice for ResNet-50 classification, Azure Custom Vision for high-level categorization, and Azure AI Vision for OCR text extraction. The results from all three services converge at the Azure OpenAI service for contextual synthesis, producing a comprehensive verification report persisted in SQL Server and returned to the frontend.

\begin{figure}[h!]
  \centering
  \begin{tikzpicture}[
      node distance=1.4cm,
      proc/.style={rectangle, draw, rounded corners=4pt,
                   minimum width=4.2cm, minimum height=0.9cm,
                   align=center, fill=blue!7, font=\small},
      ext/.style={proc, fill=orange!12},
      store/.style={proc, fill=gray!15},
      arr/.style={->, >=stealth},
      wide/.style={arr, line width=0.9pt}
    ]

    \node[ext]   (farmer)     {Farmer — Upload Seed Images};
    \node[proc, below=of farmer]  (proxy)  {Next.js Proxy Route\\ \texttt{/api/proxy/verification}};
    \node[proc, below=of proxy]   (api)    {VerificationService\\ (ASP.NET Core 8)};

    % Parallel processing stage
    \node[proc, below left=1.4cm and 2.0cm of api]  (cv)     {Azure Custom Vision\\ Package Classification};
    \node[proc, below=1.4cm of api]                  (ocr)    {Azure AI Vision\\ OCR Text Extraction};
    \node[proc, below right=1.4cm and 2.0cm of api]  (resnet) {FastAPI + ResNet-50\\ DUS Physical Parameters};

    % Convergence and synthesis
    \node[proc, below=4.8cm of api]                  (openai) {Azure OpenAI GPT\\ Natural Language Synthesis};
    \node[store, below=of openai]                     (db)     {SQL Server\\ VerificationHistory Record};
    \node[ext,   below=of db]                         (resp)   {JSON Response \u2192 Frontend};

    \draw[wide] (farmer) -- (proxy);
    \draw[wide] (proxy)  -- (api);
    \draw[wide] (api)    -- (cv);
    \draw[wide] (api)    -- (ocr);
    \draw[wide] (api)    -- (resnet);
    \draw[wide] (cv)     -- (openai);
    \draw[wide] (ocr)    -- (openai);
    \draw[wide] (resnet) -- (openai);
    \draw[wide] (openai) -- (db);
    \draw[wide] (db)     -- (resp);
  \end{tikzpicture}
  \caption{Data flow diagram showing the parallel verification pipeline from image submission through AI service orchestration to result persistence and frontend response.}
  \label{fig:dfd}
\end{figure}

% ──────────────────────────────────────
\subsection{Entity-Relationship Diagram}

Figure~\ref{fig:erd} presents the simplified entity-relationship diagram for
the core domain model persisted in Microsoft SQL Server.

\begin{figure}[h!]
  \centering
  \begin{tikzpicture}[
      entity/.style={rectangle, draw, minimum width=3.4cm,
                     minimum height=1.0cm, align=center,
                     fill=green!8, font=\small\bfseries},
      attr/.style={font=\scriptsize, align=left},
      rel/.style={->, >=stealth, thick},
      label/.style={font=\scriptsize, midway, above}
    ]

    %% Entities with adjusted positioning
    \node[entity] (user)     at (0, 0)      {User};
    \node[entity] (vh)       at (6, 1.8)    {VerificationHistory};
    \node[entity] (pc)       at (6, -1.8)   {ProductComplaint};
    \node[entity] (ca)       at (11.5, -1.8) {ComplaintAction};
    \node[entity] (brr)      at (11.5, 1.8)  {BatchRiskRegistry};
    \node[entity] (rd)       at (0, -4.5)    {ReferenceData};
    \node[entity] (rt)       at (6, -4.5)    {RefreshToken};

    %% Attributes for each entity
    \node[attr, below=0.08cm of user] {%
      \scriptsize
      \texttt{UserId (PK)}\\
      \texttt{Username, Email}\\
      \texttt{PasswordHash}\\
      \texttt{Role: enum}\\
      \texttt{State, District}
    };

    \node[attr, below=0.08cm of vh] {%
      \scriptsize
      \texttt{HistoryId (PK)}\\
      \texttt{ImageUrl, OCRText}\\
      \texttt{ClassificationResult}\\
      \texttt{ConfidenceScore}\\
      \texttt{UserId (FK)}
    };

    \node[attr, below=0.08cm of pc] {%
      \scriptsize
      \texttt{ComplaintId (PK)}\\
      \texttt{IssueType, District}\\
      \texttt{BatchNumber, Severity}\\
      \texttt{Status: enum}\\
      \texttt{UserId, OfficerId (FK)}
    };

    \node[attr, below=0.08cm of ca] {%
      \scriptsize
      \texttt{ActionId (PK)}\\
      \texttt{ComplaintId (FK)}\\
      \texttt{Description}\\
      \texttt{ActionDate}
    };

    \node[attr, below=0.08cm of brr] {%
      \scriptsize
      \texttt{BatchNumber (PK)}\\
      \texttt{TotalComplaints}\\
      \texttt{AvgSeverityScore}\\
      \texttt{RiskLevel: enum}
    };

    %% Relationships
    \draw[rel] (user) -- node[label]{1 : *} (vh);
    \draw[rel] (user) -- node[label, below]{1 : *}  (pc);
    \draw[rel] (pc)   -- node[label]{1 : *}     (ca);

  \end{tikzpicture}
  \caption{Entity-relationship diagram depicting the core domain model with primary entities, attributes, and cardinality relationships.}
  \label{fig:erd}
\end{figure}

\noindent
Key design decisions reflected in the domain model include:

\begin{itemize}
  \item \textbf{Soft Delete Pattern} — The \texttt{VerificationHistory} and \texttt{ProductComplaint} entities implement soft-delete semantics through \texttt{IsDeleted} and \texttt{DeletedAt} fields, preserving historical records for compliance auditing while enabling logical deletion without database record removal.

  \item \textbf{Third Normal Form Normalization} — Seed package information, verification records, and complaint data are stored in separate entities to eliminate data redundancy and ensure referential integrity. This normalization strategy supports efficient querying and updates while minimizing storage overhead.

  \item \textbf{Indexing Strategy} — Non-clustered indices are created on frequently-queried columns (\texttt{User.Email}, \texttt{VerificationHistory.UserId}, \texttt{ProductComplaint.Status}, \texttt{ProductComplaint.District}) to accelerate dashboard query execution and improve overall system responsiveness.

  \item \textbf{Denormalized Risk Registry} — The \texttt{BatchRiskRegistry} entity implements controlled denormalization to provide fast aggregated risk analytics for officers. A background job updates the registry whenever complaint state transitions occur, balancing consistency with query performance.
\end{itemize}

\newpage

% ──────────────────────────────────────────────────────────────
\section{Security and Deployment Considerations}
% ──────────────────────────────────────────────────────────────

\subsection{Authentication and Authorization}

The platform implements a multi-layered security architecture. Frontend authentication state is managed through Zustand stores synchronized with browser localStorage, while backend authentication is enforced through ASP.NET Core Identity and JWT tokens. Role-based access control (RBAC) policies restrict endpoint access based on user roles (Farmer, Officer, Administrator). All credentials are validated against hashed values stored in SQL Server using industry-standard algorithms.

\subsection{Data Protection}

Images uploaded for verification are encrypted and persisted in Azure Blob Storage with access tokens validated against user identity. Sensitive data (connection strings, API keys) are never hardcoded and are instead injected through environment variables. All database queries utilize parameterized queries via Entity Framework Core to prevent SQL injection attacks.

\subsection{Containerization and Deployment}

The entire application stack is containerized using Docker, enabling consistent deployment across development, staging, and production environments. The backend is deployed on Render Cloud as a containerized service; the frontend is deployed on Vercel with server-side rendering enabled; the ML microservice is deployed on GCP Cloud Run with automatic scaling based on traffic.